\chapter{Open source: contributing to maya}
\label{ch:contributing}

\begin{aside}
\maya{} is open source. Contributions large and small are
welcome. This chapter is the orientation.
\end{aside}

\section{Where the code lives}

The \code{maya/} subdirectory of the repo. Header files
under \code{maya/include/maya/}. Tests under
\code{maya/tests/}.

\section{Building}

\begin{lstlisting}
cmake -B build
cmake --build build -j10
ctest --test-dir build
\end{lstlisting}

A C++26-capable compiler is required.

\section{Issues and PRs}

\begin{itemize}[leftmargin=*]
  \item Bugs: file an issue with reproduction steps.
  \item Features: discuss in an issue before
    implementation.
  \item Pull requests: keep them focused, one concern per
    PR.
\end{itemize}

\section{Style}

British spelling. Short paragraphs in docs. Code style
matches the existing.

\section{Tests}

Every change needs a test. Snapshots for output;
unit tests for logic; integration tests for flows.

\section{Documentation}

The book is in \code{book/}. Edits welcome. Existing-
chapter expansions and new chapters both add value.

\section{Code of conduct}

Be respectful. Disagree without disagreeing. Newcomers are
welcome.

\section{Recap}

\begin{recap}
\begin{itemize}[leftmargin=*]
  \item Build with cmake; test with ctest.
  \item Issues before features.
  \item British spelling; short paragraphs.
  \item Tests required; docs welcome.
\end{itemize}
\end{recap}

\section*{Workshop}

\begin{enumerate}[leftmargin=*]
  \item Fix a small typo somewhere and open a PR.
  \item Add a test for an existing widget.
  \item Add a paragraph to a chapter in this book.
\end{enumerate}
